\section{Coexistence criterion for two-strain models with essentially rank-one siphon invasion matrices and one
input species}

\beT[Coexistence criterion for two-strain models with essentially rank-one siphon invasion matrices and one
input species]
\lbl{t:2}
Consider a two-strain model with one susceptible/input coordinate $s$,
two minimal siphons $\Sigma_1,\Sigma_2$, and siphon variables
\[
        I_k\in\RR^{m_k}_{\ge0},
        \qquad k\eq1,2.
\]
Let $E_1,E_2$ be the one-strain boundary equilibria, with susceptible
coordinates $s_1,s_2$, and let
\[
        R_2^{E_1},
        \qquad
        R_1^{E_2}
\]
denote the invasion numbers.

Assume:
\begin{enumerate}[label=\textup{(A\arabic*)},
                  align=left, leftmargin=*, labelindent=0pt]

\item\label{A:rankone}\textup{(RANK ONE)}
The strain equations may be written as
\[
        I_k'=M_k(s,z_k)I_k,
        \qquad
        M_k(s,z_k)=w_k(s,z_k)\pi_k^T-\diag(v_k),
\]
with
\[
        w_k(s,z_k)>0,
        \qquad
        \pi_k>0,
        \qquad
        v_k>0,
\]
and $M_k(s,0)$ irreducible for every $s>0$.

\item\label{A:RUR}\textup{(SCALAR RUR POLYNOMIAL)}
After imposing the rank-one degeneracy equations
\[
        \pi_k^T\diag(v_k)^{-1}w_k\eq1,
        \qquad k\eq1,2,
\]
all variables outside
\[
        \{s\}\cup\Sigma_1\cup\Sigma_2
\]
are algebraically eliminated, and the remaining fixed-point equations
reduce to a scalar polynomial
\[
        P(s)\eq0.
\]

\item\label{A:norm}\textup{(BOUNDARY SIGN CONDITIONS)}
The polynomial $P$ satisfies the following sign condition at the
boundary of the positive reconstruction interval:
\[
        P(0)\cdot P(s_j) < 0
        \qquad
        \text{whenever } R_{3-j}^{E_j}>1,\quad j\eq1,2.
\]
In the generic case $P(0)>0$ this reduces to $P(s_j)<0$ when
$R_{3-j}^{E_j}>1$, which is then implied by \textup{(A4)} below.
In the case $P(0)<0$ (e.g.\ when the coexistence polynomial has a
unique positive root with $P\to+\infty$), the same sign-change
argument applies with the inequality reversed.

\item\label{A:factor}\textup{(POSITIVE BOUNDARY FACTORIZATION)}
Direct evaluation at $s=s_1,s_2$ gives
\[
        P(s_1)
        \eq
        -K_1\bigl(R_2^{E_1}-1\bigr),
\]
\[
        P(s_2)
        \eq
        -K_2\bigl(R_1^{E_2}-1\bigr),
\]
with
\[
        K_1,K_2>0.
\]

\item\label{A:recon}\textup{(POSITIVE RECONSTRUCTION)}
Every root
\[
        s^*\in(0,\min\{s_1,s_2\})
\]
of $P(s)\eq0$ satisfies $R_k(s^*)>1$ for $k\eq1,2$, yields positive
values for all eliminated variables, and yields positive amplitudes
\[
        c_1^*,c_2^*>0
\]
in the two siphon blocks.

\end{enumerate}

Then
\[
\begin{cases}
s_1<s_2,\ R_2^{E_1}>1
&\Longrightarrow
\exists\,E_{12}\in\RR^n_{>0},
\\[1ex]
s_2<s_1,\ R_1^{E_2}>1
&\Longrightarrow
\exists\,E_{12}\in\RR^n_{>0}.
\end{cases}
\]
Moreover, the susceptible coordinate of $E_{12}$ satisfies
$s^*<\min\{s_1,s_2\}$, so both reproduction functions
$R_k(s^*)>1$, $k\eq1,2$, confirming that both siphon blocks are
active at $E_{12}$.
\eeT

\begin{proof}
Assume first $s_1<s_2$.  If $R_2^{E_1}>1$, then by \ref{A:factor},
\[
        P(s_1)\eq-K_1(R_2^{E_1}-1)<0.
\]
By \ref{A:norm}, $P(0)$ and $P(s_1)$ have opposite signs, so the
intermediate value theorem gives $s^*\in(0,s_1)$ with $P(s^*)\eq0$.
Because $s_1<s_2$,
\[
        s^*<s_1<s_2,
        \qquad\text{hence}\qquad
        s^*\in(0,\min\{s_1,s_2\}).
\]

By \ref{A:recon}, all eliminated variables are positive,
$R_k(s^*)>1$ for $k\eq1,2$, and the amplitudes satisfy
$c_1^*,c_2^*>0$.

It remains to reconstruct the siphon variables $I_k^*$.
At $s=s^*$, the rank-one degeneracy equations in \ref{A:RUR} give
\[
        \pi_k^T\diag(v_k)^{-1}w_k(s^*,z_k(s^*))\eq1,
        \qquad k\eq1,2.
\]
Therefore
\begin{align*}
        M_k(s^*,z_k(s^*))
        \diag(v_k)^{-1}w_k(s^*,z_k(s^*))
        &\eq
        \bigl[w_k(s^*,z_k(s^*))\pi_k^T-\diag(v_k)\bigr]
        \diag(v_k)^{-1}w_k(s^*,z_k(s^*))
        \\
        &\eq
        w_k(s^*,z_k(s^*))
        \underbrace{\pi_k^T\diag(v_k)^{-1}w_k(s^*,z_k(s^*))}_{\eq1}
        -w_k(s^*,z_k(s^*))
        \eq0.
\end{align*}
Thus $\diag(v_k)^{-1}w_k(s^*,z_k(s^*))\in\ker M_k(s^*,z_k(s^*))$.
Since $R_k(s^*)>1$ and $M_k(s^*,0)$ is irreducible by \ref{A:rankone},
Theorem~\ref{t:relayRF} guarantees that the Perron eigenvalue of
$M_k(s^*,0)$ is positive and simple; hence $\ker M_k(s^*,z_k(s^*))$
is one-dimensional (as $M_k(s^*,z_k(s^*))$ is a rank-one perturbation
of the nonsingular diagonal $-\diag(v_k)$ with exactly one imposed
degeneracy).  Consequently
\[
        \ker M_k(s^*,z_k(s^*))
        \eq
        \bigl\langle
        \diag(v_k)^{-1}w_k(s^*,z_k(s^*))
        \bigr\rangle.
\]
The reconstructed siphon variables are therefore
\[
        I_k^*
        \eq
        c_k^*\diag(v_k)^{-1}w_k(s^*,z_k(s^*)),
        \qquad c_k^*>0.
\]
Since $v_k>0$ and $w_k(s^*,z_k(s^*))>0$, we obtain $I_k^*>0$ for
$k\eq1,2$.  Together with the positive reconstruction of the
eliminated variables, the reconstructed fixed point belongs to
$\RR^n_{>0}$.

The case $s_2<s_1$ is symmetric.  If $R_1^{E_2}>1$, then
\ref{A:factor} gives $P(s_2)<0$, and \ref{A:norm} gives a sign change
on $(0,s_2)$, yielding $s^*\in(0,s_2)\subset(0,\min\{s_1,s_2\})$.
The same rank-one kernel calculation gives $I_k^*>0$ for $k\eq1,2$,
hence $E_{12}\in\RR^n_{>0}$.
\end{proof}

\beR[role of elimination theory in coexistence results]
The coexistence criterion uses CRNT and NGM structure to identify the
rank-one degeneracy equations, but it also requires a genuinely
elimination-theoretic input.  The scalar polynomial $P(s)$ must exist,
its boundary values must factor through the two invasion quantities via
\ref{A:factor}, and its roots in $(0,\min\{s_1,s_2\})$ must
reconstruct positive coordinates via \ref{A:recon}.  The sign
condition \ref{A:norm} replaces the naive requirement $P(0)>0$: in
models such as the OSN model with $D_0>0$, one has $P(0)<0$ and the
IVT argument runs in the opposite direction, but the conclusion is
the same.  Thus coexistence is not obtained from siphon or NGM
structure alone; it is obtained when the elimination problem is
compatible with the invasion inequalities and the sign-change condition
holds at the boundary of the positive reconstruction interval.
\eeR